\section{Experiment}\label{sec:Experiment}
    Write the Experiment parts here.%写一个阶段实验章的开场白
    \subsection{Implementation Details}\label{subsec:Implementation Details}%实验细节
        \paragraph{Datasets.} We evaluate on four widely-used change detection datasets: LEVIR-CD, WHU-CD, GZ-CD, and CLCD. These datasets cover various geographic regions, land cover types, and change patterns, providing a comprehensive testbed for cross-dataset generality.%后续自己重写数据集描述：简介，引用，

        \paragraph{Evaluation Metrics.} We report six metrics: Precision (Pre), Recall (Rec), F1-score (F1), Intersection over Union (IoU), Overall Accuracy (OA), and Kappa coefficient. These metrics capture different aspects of detection quality, including pixel-wise classification, region overlap, and agreement beyond chance.%后续自己重写指标描述：简介，公式，意义

        \paragraph{Model Details.} The model details, training settings, and other implementation specifics will be described in the next sections.%网络模型的介绍，对比方法中所用的神经网络模型，

        \paragraph{Hyperparameters.} % 主要超参的值，可以用表格的形式

        \paragraph{Hardware Details.} % 介绍服务器处理器显卡内存信息系统等硬件环境信息
    

    \subsection{Main Results}\label{subsec:Main Results}%表格呈现本方法与主要对比方法的结果同时需要针对表格有一段量化分析讨论的表述


    \subsection{Results on Pseudo-Change Datasets}\label{subsec:Results on Pseudo-Change Datasets}%表格形式在伪变化数据集上与主流方法作对比同时针对表格有一段量化分析讨论的表述


    \subsection{Ablation Study}\label{subsec:Ablation Study}%表格形式展示消融实验结果同时针对表格有一段量化分析讨论的表述


    \subsection{Visualization Analysis}\label{subsec:Visualization Analysis}%可视化结果展示同时针对可视化结果有一段分析讨论的表述,用来挑选对我有利的可视化结果，展示我方法在伪变化抑制和边界优化方面的优势


    \subsection{Model Complexity and Efficiency}\label{subsec:Model Complexity and Efficiency}%表格形式展示模型复杂度和效率结果同时针对表格有一段量化分析讨论的表述，训练的时间、显存的开销，推理的时间和显存的开销等 推理时间、参数量，python包calops等工具计算模型的复杂度和效率指标,统计模型的参数量、FLOPs、训练时间、推理时间等指标，并与对比方法进行比较，分析模型在实际应用中的效率和可行性。


    \subsection{Discussion}\label{subsec:Discussion}%对实验结果的总结和分析，强调我方法在伪变化抑制和边界优化方面的优势，同时也指出可能的不足和未来的改进方向
        

